% =====================================================================
%  M5 -- Heavy-ball momentum vs plain GD: spectral comparison
% =====================================================================
\subsection{M5: Heavy-ball momentum vs.\ plain GD}
\setprob{M5}

\begin{frame}{\probtitle{M5}{Heavy-ball momentum against plain GD}}
  \begin{probbox}[Problem M5 \hfill Mathematical \quad $\star\star\star$]
    The heavy-ball (Polyak) iteration is
    $\x_{k+1} = \x_k - \alpha\nabla f(\x_k) + \beta(\x_k - \x_{k-1})$.
    Model each eigen-direction by $f(x) = \tfrac12\lambda x^2$ with $\lambda\in[\mu,L]$.
    \begin{enumerate}[(a)]
      \item Show that $x_{k+1} = (1+\beta-\alpha\lambda)x_k - \beta x_{k-1}$, and derive its characteristic polynomial.
      \item Show that when the roots are complex, $|r_1| = |r_2| = \sqrt\beta$, \emph{whatever $\lambda$ is}.
      \item For $\mu = 1$, $L = 100$, take
            $\alpha = \frac{4}{(\sqrt L+\sqrt\mu)^2}$ and $\beta = \bigl(\frac{\sqrt L-\sqrt\mu}{\sqrt L+\sqrt\mu}\bigr)^2$.
            Compute them, and show that every $\lambda\in[1,100]$ gives complex or repeated roots.
      \item How many iterations do GD (best $\alpha = \frac{2}{\mu+L}$) and heavy ball need to cut the error by $10^{6}$?
      \item What goes on at exactly $\lambda = \mu$?
    \end{enumerate}
  \end{probbox}
\end{frame}

\begin{frame}{\soltitle{M5}{1}{4}}
  \textbf{(a)} With $\nabla f = \lambda x$:
  $\;x_{k+1} = x_k - \alpha\lambda x_k + \beta x_k - \beta x_{k-1} = (1+\beta-\alpha\lambda)\,x_k - \beta\,x_{k-1}.$

  Trying $x_k = r^k$ gives the \textbf{characteristic polynomial}
  \[
    \ans{r^2 - (1+\beta-\alpha\lambda)\,r + \beta = 0},
    \qquad\text{with companion matrix } \mathbf T = \begin{pmatrix}1+\beta-\alpha\lambda & -\beta\\ 1 & 0\end{pmatrix}.
  \]
  The iteration converges iff the spectral radius $\rho(\mathbf T) = \max|r_i|$ is below 1.

  \vspace{4pt}
  \textbf{(b)} Vieta gives $r_1 r_2 = \beta$. If the roots are complex conjugates, $r_2 = \bar r_1$, so
  \[
    |r_1|^2 = r_1\bar r_1 = r_1 r_2 = \beta \;\Longrightarrow\; \ans{|r_{1,2}| = \sqrt\beta}.
  \]
  The roots are complex exactly when the discriminant is negative:
  \[
    (1+\beta-\alpha\lambda)^2 < 4\beta
    \iff (1-\sqrt\beta)^2 < \alpha\lambda < (1+\sqrt\beta)^2 .
  \]
\end{frame}

\begin{frame}{\soltitle{M5}{2}{4}}
  \textbf{(c)} Here $\kappa = L/\mu = 100$ and $\sqrt\kappa = 10$, so
  \[
    \alpha = \frac{4}{(10+1)^2} = \ans{\frac{4}{121}\approx 0.03306},
    \qquad
    \sqrt\beta = \frac{9}{11},\;\; \beta = \ans{\frac{81}{121}\approx 0.6694}.
  \]
  The edges of the complex region are
  \[
    (1-\sqrt\beta)^2 = \Bigl(\frac{2}{11}\Bigr)^2 = \frac{4}{121} = \alpha\mu,
    \qquad
    (1+\sqrt\beta)^2 = \Bigl(\frac{20}{11}\Bigr)^2 = \frac{400}{121} = \alpha L .
  \]
  So as $\lambda$ runs over $[1,100]$, $\alpha\lambda$ covers \emph{exactly} $[\tfrac4{121},\tfrac{400}{121}]$: complex roots inside,
  a double root at each end. Therefore
  \[
    \rho(\mathbf T) = \sqrt\beta = \frac{9}{11} \approx 0.818 \quad\text{for \emph{every} } \lambda\in[1,100].
  \]
  \small\textcolor{mGrey}{That is the whole point of this tuning: it puts the entire spectrum on the circle $|r| = \sqrt\beta$.}
\end{frame}

\begin{frame}{\soltitle{M5}{3}{4}}
  \begin{columns}[T]
    \begin{column}{0.53\textwidth}
      \textbf{(d)} Contraction factors, and iterations from $k \ge \ln 10^6 / \ln(1/\rho)$:
      \begin{align*}
        \rho_{\text{GD}} &= \tfrac{\kappa-1}{\kappa+1} = \tfrac{99}{101}, &
        k_{\text{GD}} &\ge \tfrac{13.8155}{0.0200} \Rightarrow \ans{691}\\[2pt]
        \rho_{\text{HB}} &= \tfrac{\sqrt\kappa-1}{\sqrt\kappa+1} = \tfrac{9}{11}, &
        k_{\text{HB}} &\ge \tfrac{13.8155}{0.2007} \Rightarrow \ans{69}
      \end{align*}
      That is about $10 = \sqrt\kappa$ times faster. In general the cost drops from
      $O(\kappa\log\tfrac1\varepsilon)$ to $O(\sqrt\kappa\log\tfrac1\varepsilon)$.
    \end{column}
    \begin{column}{0.44\textwidth}
      \centering
      \includegraphics[width=\linewidth]{fig_hb_spectrum}
    \end{column}
  \end{columns}
\end{frame}

\begin{frame}{\soltitle{M5}{4}{4}}
  \textbf{(e) Repeated roots at the ends.} At $\lambda = \mu$ the discriminant is $0$ and there is a double root
  $r = \tfrac12(1+\beta-\alpha\mu) = \sqrt\beta = \tfrac9{11}$, so
  \[
    x_k = (c_1 + c_2\,k)\,r^k \qquad\text{(an extra factor of $k$ appears).}
  \]
  \begin{itemize}\small
    \item So the error is \emph{not} bounded by $1\cdot r^k$. From $x_0 = x_{-1} = 1$ with $\lambda = \mu$,
          heavy ball needs \alert{83} iterations for $|x_k|<10^{-6}$, not 69 (checked numerically); GD needs 691.
    \item Rates like $\rho_{\text{HB}}$ are \emph{asymptotic}: $\norm{\e_k}\le C(k)\,\rho^k$, where $C(k)$ grows at most polynomially.
  \end{itemize}
  \begin{warnbox}[Beyond quadratics, be careful]
    For a general strongly convex $f$, heavy ball with these parameters is not guaranteed to converge.
    Lessard, Recht and Packard give a counterexample~\cite{lessard2016}; Nesterov's method has a proof.
  \end{warnbox}
\end{frame}
